\clearpage
\SourceLeaf{194}
\thispagestyle{empty}
\pdfbookmark[0]{Allgemeine Lösung der Aufgabe, die Teile einer gegebenen Fläche auf eine andere abzubilden}{abbildung}
\vspace*{.35in}
\begin{center}
{\large ALLGEMEINE LÖSUNG DER AUFGABE}\par\bigskip
{\Large DIE TEILE EINER GEGEBENEN FLÄCHE}\par\bigskip
AUF EINE ANDERE GEGEBENE FLÄCHE SO ABZUBILDEN\par\bigskip
DASS DIE ABBILDUNG DEM ABGEBILDETEN\par\bigskip
IN DEN KLEINSTEN TEILEN ÄHNLICH WIRD\par\vspace{28pt}
VON\par\vspace{28pt}
{\large CARL FRIEDRICH GAUSS}\par\vspace{30pt}
{\small ALS ANTWORT AUF DIE VON DER KÖNIGLICHEN GESELLSCHAFT DER WISSENSCHAFTEN\par
IN KOPENHAGEN FÜR MDCCCXXII GESTELLTE PREISFRAGE.}\par\bigskip
\end{center}
\begin{flushright}‘Von diesen aus wird der Weg zu Größerem gebahnt.’\end{flushright}
\vfill
\begin{center}
\rule{.8\linewidth}{.3pt}\par\smallskip
Astronomische Abhandlungen herausgegeben von H. C. \textsc{Schumacher}.\par
Drittes Heft. Altona 1825.\par\smallskip
\rule{.8\linewidth}{.3pt}
\end{center}
\clearpage
\SourceLeaf{196}
\fancyhead[L]{\small ALLGEMEINE LÖSUNG DER AUFGABE}
\pdfbookmark[1]{Vorrede, Dezember 1822}{abbildung-vorrede}
Der Verfasser dieser Abhandlung glaubte, die zweimalige Wahl der Aufgabe, die ihren Gegenstand bildet, als Beweis für die Wichtigkeit ansehen zu müssen, die die Königliche Gesellschaft ihr beimisst. Das ermutigte ihn, ihr seine schon vor längerer Zeit gefundene Lösung vorzulegen; sonst hätte ihn die späte Kenntnis der Preisfrage davon abgehalten. Er bedauert, dass dieser Umstand ihn gezwungen hat, sich fast ausschließlich auf das Wesentliche und auf Hinweise zu einigen naheliegenden Anwendungen bei Kartenprojektionen und in der höheren Geodäsie zu beschränken. Ohne den nahen Abgabetermin hätte er gern einige Nebenfragen weiterentwickelt und die vielfältigen Anwendungen in der höheren Geodäsie ausführlich ausgearbeitet; dies muss er sich nun für eine andere Zeit und einen anderen Ort vorbehalten.

Im Dezember 1822.
\begin{center}\rule{.3\linewidth}{.3pt}\end{center}
\clearpage
\SourceLeaf{198}
\begin{center}
{\large ALLGEMEINE LÖSUNG DER AUFGABE}\par\bigskip
DIE TEILE EINER GEGEBENEN FLÄCHE\par\medskip
AUF EINE ANDERE GEGEBENE FLÄCHE SO ABZUBILDEN\par\medskip
DASS DIE ABBILDUNG DEM ABGEBILDETEN\par\medskip
IN DEN KLEINSTEN TEILEN ÄHNLICH WIRD.\par\bigskip
\rule{.25\linewidth}{.3pt}
\end{center}
\renewcommand{\Article}[1]{\par\medskip\Needspace{4\baselineskip}\pdfbookmark[1]{Art. #1}{abbildung-art-#1}\centerline{\textbf{#1.}}\nobreak\smallskip}
\Article{1}
Die Natur einer gekrümmten Fläche wird durch eine Gleichung zwischen den Koordinaten \(x,y,z\) jedes ihrer Punkte bestimmt. Aufgrund dieser Gleichung kann jede dieser drei veränderlichen Größen als Funktion der beiden anderen betrachtet werden. Noch allgemeiner ist es, zwei weitere veränderliche Größen \(t,u\) einzuführen und jede der Größen \(x,y,z\) als Funktion von \(t\) und \(u\) darzustellen; dabei gehören, zumindest allgemein gesprochen, bestimmte Werte von \(t\) und \(u\) stets zu einem bestimmten Punkt der Oberfläche und umgekehrt.

\Article{2}
Für eine zweite gekrümmte Fläche sollen \(X,Y,Z,T,U\) entsprechende Bedeutungen haben wie jeweils \(x,y,z,t,u\) für die erste.

\Article{3}
Die erste Fläche auf die zweite \textit{abzubilden} heißt, eine Regel festzulegen, nach der jedem Punkt der ersten Fläche ein bestimmter Punkt der zweiten entsprechen soll. Das geschieht, indem \(T\) und \(U\) bestimmten Funktionen der beiden veränderlichen Größen \(t\) und \(u\) gleichgesetzt werden.
\SourceLeaf{199}
Soll die Abbildung bestimmte Bedingungen erfüllen, so dürfen diese Funktionen nicht mehr beliebig sein. Da dadurch auch \(X,Y,Z\) zu Funktionen von \(t\) und \(u\) werden, müssen diese Funktionen neben der Bedingung, die die Natur der zweiten Fläche vorschreibt, auch die Bedingung erfüllen, die für die Abbildung gelten soll.

